\paragraph*{04.05.04.03 Vorgehensweise für Termine und Phasen, ggf. Sprints festlegen}
\kcilabel{para:04.05.04.03_VorgehensweiseFuerTermineUndPhasenUndSprints}{04.05.04.03}
\label{para:04.05.04.03_VorgehensweiseFuerTermineUndPhasenUndSprints}

% Task
\par Meine Aufgabe war es, die Vorgehensweise zur Festlegung von Terminen, Phasen und Sprints im \gls{wsr} zu definieren und mit dem \gls{s4}-Programm zu synchronisieren.

% Action
\par Nach meinem Einstieg im Januar 2025 analysierte ich die bisherige Terminplanung und Phasenstruktur, um eine Synchronisation der \gls{ricefw}-Aktivitäten mit den fachlichen \gls{s4}-Anforderungen sicherzustellen.

% Result
\par Eine mit dem \gls{s4}-Programm synchronisierte Terminplanung mit klarer Phasenstruktur wurde etabliert und ermöglichte effiziente Projektplanung mit definierten Meilensteinen.

% Reflect
\par Die initialen Phasenlängen waren zu lang für flexible Anforderungsreaktionen. Für Welle~2 schlug ich verkürzte Sprintzyklen (2-3 Wochen) vor, um Flexibilität zu erhöhen (vgl. \autoref{fig:AnregungenW2undLessonsLearnedW1}).